% Options for packages loaded elsewhere
\PassOptionsToPackage{unicode}{hyperref}
\PassOptionsToPackage{hyphens}{url}
\PassOptionsToPackage{dvipsnames,svgnames,x11names}{xcolor}
%
\documentclass[
  10pt,
  letterpaper,
  onecolumn]{article}
\usepackage{amsmath,amssymb}
\usepackage{iftex}
\ifPDFTeX
  \usepackage[T1]{fontenc}
  \usepackage[utf8]{inputenc}
  \usepackage{textcomp} % provide euro and other symbols
\else % if luatex or xetex
  % fontspec, NOT unicode-math: unicode-math makes math glyphs math-active,
  % which breaks \newunicodechar mappings in text mode.
  \usepackage{fontspec}
  \defaultfontfeatures{Scale=MatchLowercase}
  \defaultfontfeatures[\rmfamily]{Ligatures=TeX,Scale=1}
\fi
\IfFileExists{lmodern.sty}{\ifPDFTeX\usepackage{lmodern}\fi}{}
\ifPDFTeX\else
  % xetex/luatex font selection
  \setmainfont[]{Latin Modern Roman}
\fi
% Use upquote if available, for straight quotes in verbatim environments
\IfFileExists{upquote.sty}{\usepackage{upquote}}{}
\IfFileExists{microtype.sty}{% use microtype if available
  \usepackage[]{microtype}
  \UseMicrotypeSet[protrusion]{basicmath} % disable protrusion for tt fonts
}{}
\makeatletter
\@ifundefined{KOMAClassName}{% if non-KOMA class
  \IfFileExists{parskip.sty}{%
    \usepackage{parskip}
  }{% else
    \setlength{\parindent}{0pt}
    \setlength{\parskip}{6pt plus 2pt minus 1pt}}
}{% if KOMA class
  \KOMAoptions{parskip=half}}
\makeatother
\usepackage{xcolor}
\usepackage[margin=0.85in]{geometry}
\setlength{\emergencystretch}{3em} % prevent overfull lines
\providecommand{\tightlist}{%
  \setlength{\itemsep}{0pt}\setlength{\parskip}{0pt}}
\setcounter{secnumdepth}{5}
\newcommand{\notestitle}{Linear Algebra Notes}
% ---------------------------------------------------------------------------
%  Study-notes preamble (inlined by pandoc -H, so each .tex is standalone).
%  Engine: XeLaTeX.  Class: article, 10pt, one column.
% ---------------------------------------------------------------------------
\usepackage{amsmath,amssymb}
\usepackage{newunicodechar}
\usepackage{enumitem}
\usepackage{fancyhdr}
\usepackage{titlesec}
\usepackage{microtype}

% --- fonts ------------------------------------------------------------------
% Latin Modern for text/math; DejaVu Sans Mono for code (wide Unicode coverage).
\setmonofont{DejaVu Sans Mono}[Scale=MatchLowercase]

% --- Unicode the text font cannot draw -> real math glyphs ------------------
\newunicodechar{−}{\ensuremath{-}}
\newunicodechar{·}{\ensuremath{\cdot}}
\newunicodechar{×}{\ensuremath{\times}}
\newunicodechar{÷}{\ensuremath{\div}}
\newunicodechar{±}{\ensuremath{\pm}}
\newunicodechar{≤}{\ensuremath{\leq}}
\newunicodechar{≥}{\ensuremath{\geq}}
\newunicodechar{≈}{\ensuremath{\approx}}
\newunicodechar{≠}{\ensuremath{\neq}}
\newunicodechar{≅}{\ensuremath{\cong}}
\newunicodechar{≇}{\ensuremath{\ncong}}
\newunicodechar{∞}{\ensuremath{\infty}}
\newunicodechar{√}{\ensuremath{\sqrt{\;}}}
\newunicodechar{Σ}{\ensuremath{\Sigma}}
\newunicodechar{∫}{\ensuremath{\int}}
\newunicodechar{∈}{\ensuremath{\in}}
\newunicodechar{∘}{\ensuremath{\circ}}
\newunicodechar{‖}{\ensuremath{\|}}
\newunicodechar{ℓ}{\ensuremath{\ell}}
\newunicodechar{→}{\ensuremath{\rightarrow}}
\newunicodechar{←}{\ensuremath{\leftarrow}}
\newunicodechar{↔}{\ensuremath{\leftrightarrow}}
\newunicodechar{⇒}{\ensuremath{\Rightarrow}}
\newunicodechar{⇐}{\ensuremath{\Leftarrow}}
\newunicodechar{⇔}{\ensuremath{\Leftrightarrow}}
\newunicodechar{⌈}{\ensuremath{\lceil}}
\newunicodechar{⌉}{\ensuremath{\rceil}}
\newunicodechar{⌊}{\ensuremath{\lfloor}}
\newunicodechar{⌋}{\ensuremath{\rfloor}}
\newunicodechar{°}{\ensuremath{{}^{\circ}}}
\newunicodechar{✓}{\ensuremath{\checkmark}}
% greek
\newunicodechar{α}{\ensuremath{\alpha}}
\newunicodechar{β}{\ensuremath{\beta}}
\newunicodechar{γ}{\ensuremath{\gamma}}
\newunicodechar{ε}{\ensuremath{\epsilon}}
\newunicodechar{η}{\ensuremath{\eta}}
\newunicodechar{θ}{\ensuremath{\theta}}
\newunicodechar{λ}{\ensuremath{\lambda}}
\newunicodechar{μ}{\ensuremath{\mu}}
\newunicodechar{π}{\ensuremath{\pi}}
\newunicodechar{ρ}{\ensuremath{\rho}}
\newunicodechar{σ}{\ensuremath{\sigma}}
\newunicodechar{φ}{\ensuremath{\phi}}
\newunicodechar{ω}{\ensuremath{\omega}}
% super/subscripts
\newunicodechar{¹}{\ensuremath{{}^{1}}}
\newunicodechar{²}{\ensuremath{{}^{2}}}
\newunicodechar{³}{\ensuremath{{}^{3}}}
\newunicodechar{⁴}{\ensuremath{{}^{4}}}
\newunicodechar{⁵}{\ensuremath{{}^{5}}}
\newunicodechar{⁶}{\ensuremath{{}^{6}}}
\newunicodechar{⁷}{\ensuremath{{}^{7}}}
\newunicodechar{⁸}{\ensuremath{{}^{8}}}
\newunicodechar{ⁿ}{\ensuremath{{}^{n}}}
\newunicodechar{ˣ}{\ensuremath{{}^{x}}}
\newunicodechar{ᵀ}{\ensuremath{{}^{\mathsf{T}}}}
\newunicodechar{ᵇ}{\ensuremath{{}^{b}}}
\newunicodechar{⁻}{\ensuremath{{}^{-}}}
\newunicodechar{⁺}{\ensuremath{{}^{+}}}
\newunicodechar{₀}{\ensuremath{{}_{0}}}
\newunicodechar{₁}{\ensuremath{{}_{1}}}
\newunicodechar{₂}{\ensuremath{{}_{2}}}
\newunicodechar{₃}{\ensuremath{{}_{3}}}
\newunicodechar{ᵢ}{\ensuremath{{}_{i}}}
\newunicodechar{ₐ}{\ensuremath{{}_{a}}}
% accented letters
\newunicodechar{î}{\ensuremath{\hat{\imath}}}
\newunicodechar{ĵ}{\ensuremath{\hat{\jmath}}}
\newunicodechar{ŷ}{\ensuremath{\hat{y}}}
\newunicodechar{Ŷ}{\ensuremath{\hat{Y}}}
\newunicodechar{ȳ}{\ensuremath{\bar{y}}}
% private-use slots written by prep.py for combining-accent sequences
\newunicodechar{}{\ensuremath{\hat{f}}}
\newunicodechar{}{\ensuremath{\hat{\beta}}}
\newunicodechar{}{\ensuremath{\hat{\sigma}}}
\newunicodechar{}{\ensuremath{\hat{k}}}
\newunicodechar{}{\ensuremath{\hat{\mu}}}
\newunicodechar{}{\ensuremath{\bar{x}}}

% --- layout -----------------------------------------------------------------
\setlength{\parskip}{0.45em}
\setlength{\parindent}{0pt}
\setlist{itemsep=0.15em,parsep=0.15em,topsep=0.3em}

\titleformat{\section}{\normalfont\Large\bfseries}{\thesection}{0.7em}{}
\titleformat{\subsection}{\normalfont\large\bfseries}{\thesubsection}{0.6em}{}
\titleformat{\subsubsection}{\normalfont\normalsize\bfseries}{\thesubsubsection}{0.5em}{}
\titlespacing*{\section}{0pt}{2.0ex plus .6ex minus .2ex}{1.0ex plus .2ex}
\titlespacing*{\subsection}{0pt}{1.6ex plus .5ex minus .2ex}{0.7ex plus .2ex}
\titlespacing*{\subsubsection}{0pt}{1.2ex plus .4ex minus .2ex}{0.5ex plus .2ex}

\pagestyle{fancy}
\fancyhf{}
\renewcommand{\headrulewidth}{0.4pt}
\fancyhead[L]{\footnotesize\nouppercase\leftmark}
\fancyhead[R]{\footnotesize\textit{\notestitle}}
\fancyfoot[C]{\footnotesize\thepage}
\fancypagestyle{plain}{\fancyhf{}\renewcommand{\headrulewidth}{0pt}\fancyfoot[C]{\footnotesize\thepage}}

% keep display math from being torn off its paragraph too eagerly
\allowdisplaybreaks
\setlength{\abovedisplayskip}{0.6em}
\setlength{\belowdisplayskip}{0.6em}
\setlength{\emergencystretch}{3em}
\sloppy
\ifLuaTeX
  \usepackage{selnolig}  % disable illegal ligatures
\fi
\IfFileExists{bookmark.sty}{\usepackage{bookmark}}{\usepackage{hyperref}}
\IfFileExists{xurl.sty}{\usepackage{xurl}}{} % add URL line breaks if available
\urlstyle{same}
\hypersetup{
  pdftitle={Linear Algebra Notes},
  colorlinks=true,
  linkcolor={Maroon},
  filecolor={Maroon},
  citecolor={Blue},
  urlcolor={NavyBlue},
  pdfcreator={LaTeX via pandoc}}

\title{Linear Algebra Notes}
\author{}
\date{}

\begin{document}
\maketitle

{
\hypersetup{linkcolor=black}
\setcounter{tocdepth}{2}
\tableofcontents
}
A running file of linear-algebra concepts (with a short calculus
appendix), in Definition / Intuition / Example / Notes form. Written to
stand alone. Math is kept light but concrete: one-line formulas plus
worked 2×2 / 3×3 examples, not proofs.

Source: \textbf{3Blue1Brown --- Essence of Linear Algebra} (and
\emph{Essence of Calculus} for the appendix). The through-line is
geometric: matrices are \emph{transformations of space}, and everything
else reads off that picture. Cross-references to the ML and statistics
notes are marked \emph{(ML)} / \emph{(stat)}.

\begin{center}\rule{0.5\linewidth}{0.5pt}\end{center}

\hypertarget{vectors-and-spaces-3blue1brown}{%
\subsection{\texorpdfstring{Vectors and Spaces
\emph{(3Blue1Brown)}}{Vectors and Spaces (3Blue1Brown)}}\label{vectors-and-spaces-3blue1brown}}

\hypertarget{vectors}{%
\subsubsection{Vectors}\label{vectors}}

\textbf{Definition.} A vector can be read three ways: geometrically as
an \textbf{arrow} from the origin with a length and direction;
algebraically as an \textbf{ordered list of numbers} (its coordinates);
and abstractly as anything you can add and scale. In 2D,
\texttt{v\ =\ {[}3,\ 2{]}} means ``3 along x, 2 along y.''

\begin{equation*}
\mathbf{v}=\begin{bmatrix}v_1\\v_2\end{bmatrix},
\qquad
\|\mathbf{v}\|=\sqrt{v_1^2+v_2^2},
\qquad
a\mathbf{v}+b\mathbf{w}=\begin{bmatrix}av_1+bw_1\\av_2+bw_2\end{bmatrix}
\end{equation*}

\textbf{Intuition.} The arrow view and the list view are two languages
for one object, and linear algebra is largely translating between them.
Adding vectors = a tip-to-tail walk; scaling by a number =
stretching/squishing the arrow (flipping it if the number is negative).

\textbf{Example.} \texttt{{[}1,\ 2{]}\ +\ {[}3,\ 0{]}\ =\ {[}4,\ 2{]}};
\texttt{2·{[}1,\ 2{]}\ =\ {[}2,\ 4{]}};
\texttt{−1·{[}1,\ 2{]}\ =\ {[}−1,\ −2{]}} (same line, opposite way).

\textbf{Notes.} Coordinates only mean something relative to a chosen
basis. → Basis vectors and unit vectors, Linear combinations and span.

\hypertarget{basis-vectors-and-unit-vectors}{%
\subsubsection{Basis vectors and unit
vectors}\label{basis-vectors-and-unit-vectors}}

\textbf{Definition.} The \textbf{basis} of a coordinate system is a set
of vectors whose combinations produce every vector in the space. In 2D
the standard basis is the \textbf{unit vectors}
\texttt{î\ =\ {[}1,\ 0{]}} and \texttt{ĵ\ =\ {[}0,\ 1{]}} (length 1,
along the axes), and any \texttt{{[}x,\ y{]}\ =\ x·î\ +\ y·ĵ}.

\begin{equation*}
\begin{bmatrix}x\\y\end{bmatrix}=x\,\hat{\imath}+y\,\hat{\jmath},
\qquad
\hat{\mathbf{u}}=\frac{\mathbf{v}}{\|\mathbf{v}\|}
\end{equation*}

\textbf{Intuition.} Coordinates are just the scalars you multiply the
basis vectors by --- the basis is the ``ruler.'' Pick a different basis
and the same arrow gets different coordinates.

\textbf{Example.} \texttt{{[}3,\ 2{]}\ =\ 3·î\ +\ 2·ĵ}. A unit vector
has length 1: \texttt{{[}3,\ 4{]}} (length 5) normalizes to
\texttt{{[}3/5,\ 4/5{]}}.

\textbf{Notes.} Any linearly independent set that spans the space can
serve as a basis. → Linear combinations and span, Change of basis.

\hypertarget{linear-combinations-and-span}{%
\subsubsection{Linear combinations and
span}\label{linear-combinations-and-span}}

\textbf{Definition.} A \textbf{linear combination} of vectors
\texttt{v,\ w} is \texttt{a·v\ +\ b·w} for scalars \texttt{a,\ b}. Their
\textbf{span} is the set of \emph{all} such combinations --- every point
reachable by scaling and adding them.

\textbf{Intuition.} Fix the vectors and sweep every scalar: you trace
out their reachable region. In 2D, two vectors pointing in different
directions span the \emph{whole plane}; if they point the same way (one
is a scalar multiple of the other) the span collapses to a \emph{line}.

\textbf{Example.} In 3D: one vector spans a line; two independent
vectors span a plane (a ``sheet'' through the origin); three independent
vectors span \emph{all} of 3D space.

\textbf{Notes.} A span always passes through the origin (take
\texttt{a\ =\ b\ =\ 0}). → Linear dependence and independence, Rank and
column space.

\hypertarget{linear-dependence-and-independence}{%
\subsubsection{Linear dependence and
independence}\label{linear-dependence-and-independence}}

\textbf{Definition.} Vectors are \textbf{linearly dependent} if at least
one can be written as a combination of the others (it lies in their span
and adds no new direction). Otherwise they are \textbf{linearly
independent}.

\textbf{Intuition.} A dependent vector is redundant --- dropping it
doesn't shrink the span. Independent vectors each contribute a genuinely
new dimension.

\textbf{Example.} \texttt{{[}1,\ 0{]}} and \texttt{{[}2,\ 0{]}} are
dependent (\texttt{{[}2,\ 0{]}\ =\ 2·{[}1,\ 0{]}}; span is just a line).
\texttt{{[}1,\ 0{]}} and \texttt{{[}0,\ 1{]}} are independent (span the
plane).

\textbf{Notes.} A basis is a linearly independent set that spans the
space; its size is the space's dimension. → Linear combinations and
span, Rank and column space.

\begin{center}\rule{0.5\linewidth}{0.5pt}\end{center}

\hypertarget{linear-transformations-and-matrices-3blue1brown}{%
\subsection{\texorpdfstring{Linear Transformations and Matrices
\emph{(3Blue1Brown)}}{Linear Transformations and Matrices (3Blue1Brown)}}\label{linear-transformations-and-matrices-3blue1brown}}

\hypertarget{linear-transformations}{%
\subsubsection{Linear transformations}\label{linear-transformations}}

\textbf{Definition.} A transformation of space that (1) keeps every
\textbf{line straight} (grid lines stay parallel and evenly spaced) and
(2) keeps the \textbf{origin fixed}. It is fully determined by where it
sends the basis vectors \texttt{î} and \texttt{ĵ}.

\textbf{Intuition.} ``Linear'' = the grid stays a grid --- no curving,
no moving the origin. Since any vector is \texttt{x·î\ +\ y·ĵ}, once you
know where \texttt{î} and \texttt{ĵ} land you know where
\emph{everything} lands: the image of \texttt{{[}x,\ y{]}} is
\texttt{x·(new\ î)\ +\ y·(new\ ĵ)}.

\textbf{Example.} A 90° counter-clockwise rotation sends
\texttt{î\ =\ {[}1,\ 0{]}\ →\ {[}0,\ 1{]}} and
\texttt{ĵ\ =\ {[}0,\ 1{]}\ →\ {[}−1,\ 0{]}}. So
\texttt{{[}2,\ 1{]}\ →\ 2·{[}0,\ 1{]}\ +\ 1·{[}−1,\ 0{]}\ =\ {[}−1,\ 2{]}}.

\textbf{Notes.} ``Track the basis vectors'' is exactly why matrices
work. → Matrices as linear transformations.

\hypertarget{matrices-as-linear-transformations}{%
\subsubsection{Matrices as linear
transformations}\label{matrices-as-linear-transformations}}

\textbf{Definition.} A matrix stores a linear transformation: its
\textbf{columns are the landing spots of the basis vectors}. The 2×2
matrix \texttt{{[}{[}a,\ b{]},\ {[}c,\ d{]}{]}} sends
\texttt{î\ →\ {[}a,\ c{]}} and \texttt{ĵ\ →\ {[}b,\ d{]}}, and applying
it is
\texttt{{[}{[}a,b{]},{[}c,d{]}{]}·{[}x,\ y{]}\ =\ {[}a·x\ +\ b·y,\ c·x\ +\ d·y{]}}.

\begin{equation*}
\begin{bmatrix}a&b\\c&d\end{bmatrix}
\begin{bmatrix}x\\y\end{bmatrix}
=x\begin{bmatrix}a\\c\end{bmatrix}+y\begin{bmatrix}b\\d\end{bmatrix}
=\begin{bmatrix}ax+by\\cx+dy\end{bmatrix}
\end{equation*}

\textbf{Intuition.} Matrix-times-vector is just ``scale each landed
basis vector by the matching input coordinate, then add'' --- reading
the columns as the new \texttt{î} and \texttt{ĵ}.

\textbf{Example.} The rotation matrix
\texttt{{[}{[}0,\ −1{]},\ {[}1,\ 0{]}{]}} times \texttt{{[}2,\ 1{]}} =
\texttt{{[}0·2\ +\ (−1)·1,\ 1·2\ +\ 0·1{]}\ =\ {[}−1,\ 2{]}} ---
matching the rotation above.

\textbf{Notes.} A matrix whose columns are linearly dependent squishes
space onto a line or point (determinant 0). → Determinant, Matrix
multiplication (composition).

\hypertarget{matrix-multiplication-composition}{%
\subsubsection{Matrix multiplication
(composition)}\label{matrix-multiplication-composition}}

\textbf{Definition.} The product \texttt{M₂·M₁} is the single
transformation that does \texttt{M₁} \textbf{then} \texttt{M₂} (read
right-to-left, like function composition). Compute it by applying
\texttt{M₂} to each column of \texttt{M₁}.

\begin{equation*}
(M_2M_1)\mathbf{v}=M_2(M_1\mathbf{v}),
\qquad
(AB)_{ij}=\sum_{k}A_{ik}B_{kj},
\qquad
AB\ne BA \text{ in general}
\end{equation*}

\textbf{Intuition.} Multiplying matrices = chaining transformations.
Order matters (\texttt{M₂·M₁\ ≠\ M₁·M₂} in general) because ``rotate
then shear'' isn't ``shear then rotate.''

\textbf{Example.}
\texttt{{[}{[}0,−1{]},{[}1,0{]}{]}\ ·\ {[}{[}1,1{]},{[}0,1{]}{]}}
(rotate ∘ shear): apply the rotation to the shear's columns ---
\texttt{{[}1,0{]}\ →\ {[}0,1{]}} and \texttt{{[}1,1{]}\ →\ {[}−1,1{]}}
--- giving \texttt{{[}{[}0,−1{]},{[}1,1{]}{]}}.

\textbf{Notes.} Associative (\texttt{(AB)C\ =\ A(BC)}) but not
commutative. → Matrices as linear transformations.

\hypertarget{determinant}{%
\subsubsection{Determinant}\label{determinant}}

\textbf{Definition.} The \textbf{determinant} measures how much a
transformation \textbf{scales areas} (2D) or \textbf{volumes} (3D). For
2×2, \texttt{det{[}{[}a,b{]},{[}c,d{]}{]}\ =\ ad\ −\ bc}. For 3×3
\texttt{{[}{[}a,b,c{]},{[}d,e,f{]},{[}g,h,i{]}{]}}, expand along the top
row: \texttt{det\ =\ a(ei\ −\ fh)\ −\ b(di\ −\ fg)\ +\ c(dh\ −\ eg)}.

\begin{equation*}
\det\begin{bmatrix}a&b\\c&d\end{bmatrix}=ad-bc,
\qquad
\det\begin{bmatrix}a&b&c\\d&e&f\\g&h&i\end{bmatrix}=a(ei-fh)-b(di-fg)+c(dh-eg)
\end{equation*}

\begin{equation*}
\det(AB)=\det(A)\det(B),
\qquad
\det(A)=0 \iff A \text{ collapses space (not invertible)}
\end{equation*}

\textbf{Intuition.} A unit square (area 1) becomes a parallelogram of
area \texttt{\textbar{}det\textbar{}}. \texttt{det\ =\ 0} means space is
squished into a lower dimension (columns are dependent). A
\textbf{negative} determinant means orientation flipped --- space turned
inside-out.

\textbf{Example (2×2).}
\texttt{det{[}{[}3,1{]},{[}0,2{]}{]}\ =\ 3·2\ −\ 1·0\ =\ 6} → areas
scale ×6. \textbf{Example (3×3).}
\texttt{det{[}{[}1,2,3{]},{[}0,1,4{]},{[}5,6,0{]}{]}\ =\ 1(1·0\ −\ 4·6)\ −\ 2(0·0\ −\ 4·5)\ +\ 3(0·6\ −\ 1·5)\ =\ −24\ +\ 40\ −\ 15\ =\ 1}.

\textbf{Notes.} \texttt{det\ =\ 0} ⇔ no inverse ⇔ columns dependent ⇔
rank drops. → Systems of linear equations and the matrix inverse, Rank
and column space, Cross product.

\hypertarget{systems-of-linear-equations-and-the-matrix-inverse}{%
\subsubsection{Systems of linear equations and the matrix
inverse}\label{systems-of-linear-equations-and-the-matrix-inverse}}

\textbf{Definition.} A linear system is \texttt{A·x\ =\ b} for unknown
\texttt{x}. If \texttt{A} is invertible (\texttt{det\ A\ ≠\ 0}), the
unique solution is \texttt{x\ =\ A⁻¹·b}. For 2×2
\texttt{A\ =\ {[}{[}a,b{]},{[}c,d{]}{]}},
\texttt{A⁻¹\ =\ (1/det)·{[}{[}d,\ −b{]},\ {[}−c,\ a{]}{]}}. For 3×3,
\texttt{A⁻¹\ =\ (1/det)·adjugate} (transpose of the cofactor matrix) ---
or just solve by Gaussian elimination.

\begin{equation*}
A\mathbf{x}=\mathbf{b}
\quad\Longrightarrow\quad
\mathbf{x}=A^{-1}\mathbf{b}\ \ (\text{if }\det A\ne 0),
\qquad
\begin{bmatrix}a&b\\c&d\end{bmatrix}^{-1}=\frac{1}{ad-bc}\begin{bmatrix}d&-b\\-c&a\end{bmatrix}
\end{equation*}

\textbf{Intuition.} \texttt{A⁻¹} is the transformation that
\emph{undoes} \texttt{A} (rotate back, un-stretch). Solving
\texttt{A·x\ =\ b} asks ``which input lands on \texttt{b}?'' --- apply
the reverse transformation to \texttt{b}.

\textbf{Example (2×2).} \texttt{A\ =\ {[}{[}3,1{]},{[}0,2{]}{]}},
\texttt{det\ =\ 6}, so \texttt{A⁻¹\ =\ (1/6){[}{[}2,−1{]},{[}0,3{]}{]}}.
To solve \texttt{A·x\ =\ {[}5,\ 4{]}}:
\texttt{x\ =\ A⁻¹{[}5,4{]}\ =\ (1/6){[}2·5\ −\ 1·4,\ 3·4{]}\ =\ (1/6){[}6,\ 12{]}\ =\ {[}1,\ 2{]}}
(check: \texttt{A·{[}1,2{]}\ =\ {[}5,4{]}} ✓).

\textbf{Notes.} If \texttt{det\ A\ =\ 0} there is no inverse, and the
system has either no solution or infinitely many. → Determinant, Null
space (kernel), Cramer's rule.

\hypertarget{rank-and-column-space}{%
\subsubsection{Rank and column space}\label{rank-and-column-space}}

\textbf{Definition.} The \textbf{column space} is the span of a matrix's
columns (all outputs \texttt{A·x}). The \textbf{rank} is the
\emph{dimension} of that column space --- how many independent
directions survive the transformation.

\begin{equation*}
\operatorname{rank}(A)+\dim\bigl(\operatorname{null}(A)\bigr)=n
\qquad\text{(rank--nullity, }n=\text{number of columns)}
\end{equation*}

\textbf{Intuition.} Rank = ``how many dimensions come out the other
side.'' A 2×2 matrix has rank 2 (full rank, output fills the plane),
rank 1 (output collapses to a line), or rank 0 (everything → origin).
Full rank ⇔ invertible ⇔ \texttt{det\ ≠\ 0}.

\textbf{Example.} \texttt{{[}{[}1,2{]},{[}2,4{]}{]}} has rank 1 (column
2 = 2× column 1, so the output is a line).
\texttt{{[}{[}1,0{]},{[}0,1{]}{]}} has rank 2.

\textbf{Notes.} Rank drop ⇔ determinant 0 ⇔ a non-trivial null space
appears. → Null space (kernel), Determinant.

\hypertarget{null-space-kernel}{%
\subsubsection{Null space (kernel)}\label{null-space-kernel}}

\textbf{Definition.} The \textbf{null space} (kernel) of \texttt{A} is
every vector \texttt{x} with \texttt{A·x\ =\ 0} --- all the vectors the
transformation squishes onto the origin.

\textbf{Intuition.} When \texttt{A} collapses a dimension (rank drops),
a whole line or plane of vectors gets crushed to zero; that set is the
null space. For the homogeneous system \texttt{A·x\ =\ 0}, the null
space \textbf{is} the complete set of solutions.

\textbf{Example.} \texttt{A\ =\ {[}{[}1,2{]},{[}2,4{]}{]}} (rank 1):
\texttt{A·{[}2,−1{]}\ =\ {[}1·2\ +\ 2·(−1),\ 2·2\ +\ 4·(−1){]}\ =\ {[}0,\ 0{]}},
so \texttt{{[}2,−1{]}} and all its multiples form the null space --- a
line.

\textbf{Notes.} Full-rank (invertible) matrices have only the trivial
null space \texttt{\{0\}}. → Rank and column space, Systems of linear
equations and the matrix inverse.

\begin{center}\rule{0.5\linewidth}{0.5pt}\end{center}

\hypertarget{products-bases-and-eigen-things-3blue1brown}{%
\subsection{\texorpdfstring{Products, Bases, and Eigen-things
\emph{(3Blue1Brown)}}{Products, Bases, and Eigen-things (3Blue1Brown)}}\label{products-bases-and-eigen-things-3blue1brown}}

\hypertarget{dot-product}{%
\subsubsection{Dot product}\label{dot-product}}

\textbf{Definition.} For vectors \texttt{a,\ b}, the \textbf{dot
product} is
\texttt{a·b\ =\ Σ\ aᵢbᵢ\ =\ \textbar{}a\textbar{}\textbar{}b\textbar{}cos\ θ},
where \texttt{θ} is the angle between them. The result is a single
number (a scalar).

\begin{equation*}
\mathbf{v}\cdot\mathbf{w}=\sum_i v_iw_i=\|\mathbf{v}\|\,\|\mathbf{w}\|\cos\theta
\qquad
(\;=0 \iff \text{perpendicular})
\end{equation*}

\textbf{Intuition.} It measures \textbf{alignment}: positive when the
vectors point the same general way, zero when perpendicular, negative
when opposing. Geometrically it's the length of \texttt{a}'s projection
onto \texttt{b}, scaled by \texttt{\textbar{}b\textbar{}}.

\textbf{Example.} \texttt{{[}3,4{]}·{[}2,1{]}\ =\ 3·2\ +\ 4·1\ =\ 10};
\texttt{{[}1,0{]}·{[}0,1{]}\ =\ 0} (perpendicular).

\textbf{Notes.} The workhorse of ML: it's the operation inside every
matrix multiply, cosine similarity, and attention score. → Matrix
multiplication (composition), Self-attention (ML).

\hypertarget{cross-product}{%
\subsubsection{Cross product}\label{cross-product}}

\textbf{Definition.} In 3D, \texttt{a\ ×\ b} is a \textbf{vector}
perpendicular to both, whose length
\texttt{\textbar{}a\textbar{}\textbar{}b\textbar{}sin\ θ} equals the
\textbf{area of the parallelogram} they span. Compute it as a symbolic
determinant
\texttt{det{[}{[}î,\ ĵ,\ {]},\ {[}a₁,\ a₂,\ a₃{]},\ {[}b₁,\ b₂,\ b₃{]}{]}}.

\begin{equation*}
\mathbf{v}\times\mathbf{w}=
\det\begin{bmatrix}\hat{\imath}&\hat{\jmath}&\hat{k}\\v_1&v_2&v_3\\w_1&w_2&w_3\end{bmatrix},
\qquad
\|\mathbf{v}\times\mathbf{w}\|=\|\mathbf{v}\|\,\|\mathbf{w}\|\sin\theta
\end{equation*}

\textbf{Intuition.} Length encodes the parallelogram's area; direction
(given by the \textbf{right-hand rule}) encodes orientation.
\textbf{Order matters}: \texttt{a\ ×\ b\ =\ −(b\ ×\ a)}
(anti-commutative) --- swapping the inputs flips the result.

\textbf{Example.} \texttt{{[}1,0,0{]}\ ×\ {[}0,1,0{]}\ =\ {[}0,0,1{]}}
(x cross y points along +z by the right-hand rule). In 2D the analogous
signed area is
\texttt{a₁b₂\ −\ a₂b₁\ =\ det{[}{[}a₁,b₁{]},{[}a₂,b₂{]}{]}}.

\textbf{Notes.} The \texttt{î,\ ĵ,\ } determinant is a mnemonic for the
formula, not a literal determinant. → Determinant.

\hypertarget{cramers-rule}{%
\subsubsection{Cramer's rule}\label{cramers-rule}}

\textbf{Definition.} Solves a square system \texttt{A·x\ =\ b}
(\texttt{det\ A\ ≠\ 0}) one coordinate at a time with determinants:
\texttt{xᵢ\ =\ det(Aᵢ)\ /\ det(A)}, where \texttt{Aᵢ} is \texttt{A} with
its \texttt{i}-th column replaced by \texttt{b}.

\begin{equation*}
x_i=\frac{\det(A_i)}{\det(A)}
\qquad A_i=A \text{ with column } i \text{ replaced by } \mathbf{b}
\end{equation*}

\textbf{Intuition.} Swapping in \texttt{b} and taking a determinant
ratio isolates each coordinate through signed-area/volume scaling.
Elegant, but slow --- Gaussian elimination is far faster for real
systems.

\textbf{Example.} \texttt{A\ =\ {[}{[}2,1{]},{[}1,3{]}{]}},
\texttt{b\ =\ {[}3,5{]}}, \texttt{det\ A\ =\ 5}.
\texttt{x₁\ =\ det{[}{[}3,1{]},{[}5,3{]}{]}/5\ =\ (9−5)/5\ =\ 0.8};
\texttt{x₂\ =\ det{[}{[}2,3{]},{[}1,5{]}{]}/5\ =\ (10−3)/5\ =\ 1.4}.

\textbf{Notes.} Mainly of geometric/theoretical value; impractical for
large \texttt{n}. → Determinant, Systems of linear equations and the
matrix inverse.

\hypertarget{change-of-basis}{%
\subsubsection{Change of basis}\label{change-of-basis}}

\textbf{Definition.} Translating a vector's coordinates from one basis
to another. If \texttt{B}'s columns are an alternate basis written in
standard coordinates, then \texttt{B·{[}v{]}\_B\ =\ {[}v{]}\_standard},
and \texttt{B⁻¹·{[}v{]}\_standard\ =\ {[}v{]}\_B}. A transformation
\texttt{A} becomes \texttt{B⁻¹\ A\ B} when expressed in the new basis.

\begin{equation*}
\mathbf{v}_{\text{ours}}=A\,\mathbf{v}_{\text{theirs}},
\qquad
M_{\text{theirs}}=A^{-1}MA
\qquad\text{(conjugation = "translate, act, translate back")}
\end{equation*}

\textbf{Intuition.} The same arrow has different coordinate lists
depending on the ``language'' (basis) describing it. \texttt{B}
translates \emph{from} the alternate language \emph{to} standard;
\texttt{B⁻¹} translates back. The \texttt{B⁻¹AB} sandwich means
``translate in, apply the transformation, translate out.''

\textbf{Example.} With alternate basis \texttt{{[}2,1{]}} and
\texttt{{[}−1,1{]}} (the columns of \texttt{B}), a vector with alternate
coordinates \texttt{{[}1,1{]}} is
\texttt{B·{[}1,1{]}\ =\ {[}2\ −\ 1,\ 1\ +\ 1{]}\ =\ {[}1,\ 2{]}} in
standard coordinates.

\textbf{Notes.} Choosing an \emph{eigenbasis} makes a transformation
diagonal (pure scaling) --- the cleanest possible coordinates. →
Eigenvectors and eigenvalues, Basis vectors and unit vectors.

\hypertarget{eigenvectors-and-eigenvalues}{%
\subsubsection{Eigenvectors and
eigenvalues}\label{eigenvectors-and-eigenvalues}}

\textbf{Definition.} An \textbf{eigenvector} of \texttt{A} is a nonzero
vector that stays on its own span under the transformation --- it only
gets \textbf{stretched or squished}, never knocked off its line:
\texttt{A·v\ =\ λ·v}. The scalar \texttt{λ} is its \textbf{eigenvalue}
(the stretch factor). Find them by solving \texttt{det(A\ −\ λI)\ =\ 0}
for the eigenvalues, then \texttt{(A\ −\ λI)v\ =\ 0} for each
eigenvector.

\begin{equation*}
A\mathbf{v}=\lambda\mathbf{v}
\qquad\Longleftrightarrow\qquad
\det(A-\lambda I)=0
\qquad\text{(characteristic polynomial)}
\end{equation*}

\textbf{Intuition.} Most vectors get rotated off their span by a
transformation; eigenvectors are the special axes that merely scale,
revealing the transformation's ``natural axes.''
\texttt{det(A\ −\ λI)\ =\ 0} is required because \texttt{(A\ −\ λI)}
must squish \texttt{v} to zero (be non-invertible) for a nonzero
solution to exist.

\textbf{Example (2×2).} \texttt{A\ =\ {[}{[}2,1{]},{[}0,3{]}{]}}.
\texttt{det(A\ −\ λI)\ =\ (2−λ)(3−λ)\ =\ 0} → \texttt{λ\ =\ 2} or
\texttt{3}. For \texttt{λ\ =\ 2}:
\texttt{(A\ −\ 2I)\ =\ {[}{[}0,1{]},{[}0,1{]}{]}} ⇒ \texttt{y\ =\ 0} ⇒
\texttt{v\ =\ {[}1,\ 0{]}}. For \texttt{λ\ =\ 3}:
\texttt{(A\ −\ 3I)\ =\ {[}{[}−1,1{]},{[}0,0{]}{]}} ⇒ \texttt{y\ =\ x} ⇒
\texttt{v\ =\ {[}1,\ 1{]}}. So \texttt{A} stretches \texttt{{[}1,0{]}}
by 2 and \texttt{{[}1,1{]}} by 3.

\textbf{Notes.} A 3×3 gives a cubic in \texttt{λ} (up to three
eigenvalues), same recipe. A full basis of eigenvectors
\emph{diagonalizes} \texttt{A} (→ change of basis), making powers
\texttt{Aⁿ} trivial. In ML they underpin PCA (the top eigenvectors of
the covariance matrix). → Change of basis, Determinant, Principal
components analysis (PCA) (ML).

\begin{center}\rule{0.5\linewidth}{0.5pt}\end{center}

\hypertarget{calculus-foundations-3blue1brown}{%
\subsection{\texorpdfstring{Calculus foundations
\emph{(3Blue1Brown)}}{Calculus foundations (3Blue1Brown)}}\label{calculus-foundations-3blue1brown}}

\hypertarget{derivatives}{%
\subsubsection{Derivatives}\label{derivatives}}

\textbf{Definition.} The \textbf{derivative}
\texttt{f\textquotesingle{}(x)} is the instantaneous rate of change ---
the slope of the tangent line:
\texttt{f\textquotesingle{}(x)\ =\ lim\_\{h→0\}\ (f(x+h)\ −\ f(x))\ /\ h}.
Core rules: power \texttt{d/dx\ xⁿ\ =\ n·xⁿ⁻¹}; sum (term by term);
product
\texttt{(fg)\textquotesingle{}\ =\ f\textquotesingle{}g\ +\ fg\textquotesingle{}};
chain
\texttt{(f(g(x)))\textquotesingle{}\ =\ f\textquotesingle{}(g(x))·g\textquotesingle{}(x)}.

\begin{equation*}
\frac{df}{dx}=\lim_{h\to 0}\frac{f(x+h)-f(x)}{h},
\qquad
(fg)'=f'g+fg',
\qquad
\bigl(f(g(x))\bigr)'=f'(g(x))\,g'(x)
\end{equation*}

\textbf{Intuition.} Zoom into a curve until it looks straight; the
derivative is that local slope --- how fast the output changes per tiny
nudge of the input. The \textbf{chain rule} (a nudge propagating through
composed functions) is exactly what backpropagation runs.

\textbf{Example.} \texttt{f(x)\ =\ 3x²} →
\texttt{f\textquotesingle{}(x)\ =\ 6x}, so at \texttt{x\ =\ 2} the curve
rises 12 per unit of \texttt{x}. \texttt{d/dx\ sin(x²)\ =\ cos(x²)·2x}
(chain rule).

\textbf{Notes.} Handy standards: \texttt{d/dx\ eˣ\ =\ eˣ},
\texttt{d/dx\ ln\ x\ =\ 1/x}, \texttt{d/dx\ sin\ x\ =\ cos\ x}. → Taylor
series, Backpropagation (ML), Gradient descent (ML).

\hypertarget{integrals-and-area-under-the-curve}{%
\subsubsection{Integrals and area under the
curve}\label{integrals-and-area-under-the-curve}}

\textbf{Definition.} The \textbf{definite integral}
\texttt{∫ₐᵇ\ f(x)\ dx} is the (signed) \textbf{area under the curve}
from \texttt{a} to \texttt{b}. The \textbf{Fundamental Theorem of
Calculus}: if \texttt{F\textquotesingle{}\ =\ f}, then
\texttt{∫ₐᵇ\ f\ dx\ =\ F(b)\ −\ F(a)} --- integration undoes
differentiation.

\begin{equation*}
\int_a^b f(x)\,dx=\lim_{\Delta x\to0}\sum_i f(x_i)\Delta x,
\qquad
\frac{d}{dx}\int_a^x f(t)\,dt=f(x)
\quad\text{(fundamental theorem)}
\end{equation*}

\textbf{Intuition.} Chop the region into infinitely many thin rectangles
of width \texttt{dx} and add them up. Finding an antiderivative
\texttt{F} (a function whose slope is \texttt{f}) gives that area
exactly, with no summing.

\textbf{Example.} \texttt{∫₀²\ 3x²\ dx\ =\ {[}x³{]}₀²\ =\ 8\ −\ 0\ =\ 8}
(since \texttt{d/dx\ x³\ =\ 3x²}). Area below the axis counts as
negative.

\textbf{Notes.} Basic antiderivatives: \texttt{∫\ xⁿ\ dx\ =\ xⁿ⁺¹/(n+1)}
(n ≠ −1), \texttt{∫\ 1/x\ dx\ =\ ln\textbar{}x\textbar{}},
\texttt{∫\ eˣ\ dx\ =\ eˣ}. Underlies probability densities, where area =
probability. → Derivatives, Probability distributions (stat).

\hypertarget{taylor-series}{%
\subsubsection{Taylor series}\label{taylor-series}}

\textbf{Definition.} Approximates a smooth function near a point
\texttt{a} by a polynomial built from its derivatives:
\texttt{f(x)\ ≈\ f(a)\ +\ f\textquotesingle{}(a)(x−a)\ +\ f\textquotesingle{}\textquotesingle{}(a)/2!·(x−a)²\ +\ f\textquotesingle{}\textquotesingle{}\textquotesingle{}(a)/3!·(x−a)³\ +\ …}.

\begin{equation*}
f(x)=\sum_{n=0}^{\infty}\frac{f^{(n)}(a)}{n!}(x-a)^n
\qquad
e^x=1+x+\frac{x^2}{2!}+\frac{x^3}{3!}+\dots
\end{equation*}

\textbf{Intuition.} Match the function's value, then its slope, then its
curvature, then higher bends --- each term corrects the previous fit, so
more terms means a better local approximation. Near \texttt{a}, a
handful of terms is often plenty.

\textbf{Example.} Around \texttt{a\ =\ 0}:
\texttt{eˣ\ ≈\ 1\ +\ x\ +\ x²/2\ +\ x³/6\ +\ …} and
\texttt{cos\ x\ ≈\ 1\ −\ x²/2\ +\ x⁴/24\ −\ …}. At \texttt{x\ =\ 0.1},
\texttt{e\^{}0.1\ ≈\ 1\ +\ 0.1\ +\ 0.005\ =\ 1.105} (true value 1.1052).

\textbf{Notes.} The \texttt{n!} denominators fall out of repeatedly
differentiating the polynomial. The first-order Taylor term is how
optimization ``linearizes'' a function locally. → Derivatives.

\begin{center}\rule{0.5\linewidth}{0.5pt}\end{center}

\hypertarget{glossary}{%
\subsection{Glossary}\label{glossary}}

\begin{itemize}
\tightlist
\item
  \textbf{Basis} --- independent vectors whose combinations span the
  space; the coordinate ``ruler.'' →
  \protect\hyperlink{basis-vectors-and-unit-vectors}{Basis vectors and
  unit vectors}.
\item
  \textbf{Change of basis} --- translate coordinates/transformations
  between bases via \texttt{B} and \texttt{B⁻¹}. →
  \protect\hyperlink{change-of-basis}{Change of basis}.
\item
  \textbf{Column space} --- the span of a matrix's columns (all its
  outputs). → \protect\hyperlink{rank-and-column-space}{Rank and column
  space}.
\item
  \textbf{Cramer's rule} --- solve \texttt{Ax=b} via ratios of
  determinants. → \protect\hyperlink{cramers-rule}{Cramer's rule}.
\item
  \textbf{Cross product} --- 3D perpendicular vector; length =
  parallelogram area; right-hand rule. →
  \protect\hyperlink{cross-product}{Cross product}.
\item
  \textbf{Derivative} --- instantaneous slope;
  \texttt{lim\ (f(x+h)−f(x))/h}. →
  \protect\hyperlink{derivatives}{Derivatives}.
\item
  \textbf{Determinant} --- area/volume scaling factor of a
  transformation; 0 ⇒ singular. →
  \protect\hyperlink{determinant}{Determinant}.
\item
  \textbf{Dot product} ---
  \texttt{Σ\ aᵢbᵢ\ =\ \textbar{}a\textbar{}\textbar{}b\textbar{}cos\ θ};
  measures alignment. → \protect\hyperlink{dot-product}{Dot product}.
\item
  \textbf{Eigenvector / eigenvalue} --- vector that only scales under
  \texttt{A} (\texttt{Av\ =\ λv}); \texttt{λ} is the factor. →
  \protect\hyperlink{eigenvectors-and-eigenvalues}{Eigenvectors and
  eigenvalues}.
\item
  \textbf{Integral} --- signed area under a curve; the inverse of the
  derivative. →
  \protect\hyperlink{integrals-and-area-under-the-curve}{Integrals and
  area under the curve}.
\item
  \textbf{Linear combination} --- \texttt{a·v\ +\ b·w} of vectors. →
  \protect\hyperlink{linear-combinations-and-span}{Linear combinations
  and span}.
\item
  \textbf{Linear dependence / independence} --- redundant vs
  genuinely-new directions. →
  \protect\hyperlink{linear-dependence-and-independence}{Linear
  dependence and independence}.
\item
  \textbf{Linear transformation} --- keeps lines straight and the origin
  fixed; set by where the basis lands. →
  \protect\hyperlink{linear-transformations}{Linear transformations}.
\item
  \textbf{Matrix} --- stores a linear transformation; its columns are
  the images of the basis vectors. →
  \protect\hyperlink{matrices-as-linear-transformations}{Matrices as
  linear transformations}.
\item
  \textbf{Matrix inverse} --- undoes a transformation;
  \texttt{x\ =\ A⁻¹b}; exists iff \texttt{det\ ≠\ 0}. →
  \protect\hyperlink{systems-of-linear-equations-and-the-matrix-inverse}{Systems
  of linear equations and the matrix inverse}.
\item
  \textbf{Matrix multiplication} --- composition of transformations
  (applied right-to-left). →
  \protect\hyperlink{matrix-multiplication-composition}{Matrix
  multiplication (composition)}.
\item
  \textbf{Null space (kernel)} --- vectors mapped to 0; the solutions of
  \texttt{Ax=0}. → \protect\hyperlink{null-space-kernel}{Null space
  (kernel)}.
\item
  \textbf{Rank} --- dimension of the column space (surviving output
  directions). → \protect\hyperlink{rank-and-column-space}{Rank and
  column space}.
\item
  \textbf{Span} --- all linear combinations of a set of vectors. →
  \protect\hyperlink{linear-combinations-and-span}{Linear combinations
  and span}.
\item
  \textbf{Taylor series} --- polynomial approximation of a function from
  its derivatives. → \protect\hyperlink{taylor-series}{Taylor series}.
\item
  \textbf{Unit vector} --- a length-1 vector; the standard basis is
  \texttt{î,\ ĵ}. →
  \protect\hyperlink{basis-vectors-and-unit-vectors}{Basis vectors and
  unit vectors}.
\item
  \textbf{Vector} --- an arrow (length + direction) / an ordered list of
  numbers. → \protect\hyperlink{vectors}{Vectors}.
\end{itemize}

\end{document}
